\documentclass[11pt]{article}

\usepackage[review]{acl}

\usepackage{times}
\usepackage{latexsym}
\usepackage[T1]{fontenc}
\usepackage[utf8]{inputenc}
\usepackage{microtype}
\usepackage{inconsolata}
\usepackage{graphicx}
\usepackage{booktabs}
\usepackage{multirow}
\usepackage{amsmath}
\usepackage{amssymb}
\usepackage{xcolor}
\usepackage{hyperref}
\usepackage{listings}

\lstset{
  basicstyle=\ttfamily\small,
  breaklines=true,
  frame=single,
  columns=fullflexible
}

\title{Know-Surgery: A Unified Toolkit for Controllable Knowledge Update in Large Language Models}

\author{
  Mingxuan Liu \\
  Affiliation \\
  \texttt{email@example.com}
}

\begin{document}
\maketitle

\begin{abstract}
Deployed large language models (LLMs) require continuous knowledge updates to correct errors, remove harmful content, and incorporate new information.
Existing tools address these needs in isolation: knowledge editing frameworks like EasyEdit focus solely on modifying factual associations, while unlearning toolkits like OpenUnlearning target only data removal.
We present \textbf{Know-Surgery}, a unified toolkit that supports three complementary knowledge operations---\textit{injection} (adding new knowledge via parameter-efficient fine-tuning), \textit{deletion} (removing knowledge via machine unlearning), and \textit{editing} (modifying specific knowledge via locate-and-edit methods)---within a single Hydra-configured framework.
Know-Surgery integrates 48 methods across the three operations and their multimodal variants, covers five knowledge types (factual, event, conceptual, procedural, and rule-based), and extends to multimodal scenarios.
Users can add new models, methods, and datasets in three steps through a registry-based architecture.
Experiments on standard benchmarks (CounterFact, ZsRE, TOFU, MUSE) demonstrate that Know-Surgery reproduces published results while enabling direct cross-paradigm comparisons.
The toolkit is available at \url{https://github.com/xxx/know-surgery} with a Flask\,+\,Vue web interface for interactive exploration.
\end{abstract}

\section{Introduction}
\label{sec:intro}

Large language models (LLMs) inevitably accumulate outdated facts, biased associations, and sensitive information during pre-training.
Post-deployment, practitioners need to \textit{update} this knowledge without costly full retraining.
Three complementary paradigms have emerged to address different facets of this need:

\begin{itemize}
\setlength\itemsep{0pt}
\item \textbf{Knowledge Editing} modifies specific factual associations through targeted parameter updates~\cite{meng2022rome,meng2023memit,mitchell2022mend}.
\item \textbf{Machine Unlearning} removes the influence of particular training data to comply with privacy regulations or safety requirements~\cite{maini2024tofu,shi2024muse}.
\item \textbf{Knowledge Injection} adds new capabilities or domain knowledge via parameter-efficient fine-tuning~\cite{hu2022lora,liu2024dora}.
\end{itemize}

However, existing toolkits treat these paradigms in isolation.
EasyEdit~\cite{wang2024easyedit} provides a comprehensive knowledge editing framework but does not support unlearning or injection.
OpenUnlearning~\cite{dorna2025openunlearning} unifies unlearning benchmarks but lacks editing capabilities.
LLaMA-Factory~\cite{zheng2024llamafactory} focuses on fine-tuning without evaluating the impact on existing knowledge.
This fragmentation forces researchers to maintain separate codebases, making cross-paradigm comparisons difficult.

We present \textbf{Know-Surgery}, a unified toolkit that integrates all three paradigms under a single framework.
Our key contributions are:

\begin{enumerate}
\setlength\itemsep{0pt}
\item A \textbf{unified architecture} based on Hydra configuration and a registry pattern, supporting 48 methods across editing (15), unlearning (14), and injection (6)---plus 13 multimodal variants---through a single entry point.
\item Coverage of \textbf{five knowledge types}---factual, event, conceptual, procedural, and rule-based---with 8+ editing datasets and standardized evaluation.
\item \textbf{Multimodal extensibility} through a dedicated sidecar supporting vision-language model unlearning on CLEAR and MLLMU-Bench benchmarks.
\item A \textbf{three-step extension protocol} enabling users to add new models, methods, or datasets by implementing one class, registering it, and creating a YAML config.
\end{enumerate}

\section{System Overview}
\label{sec:system}

Know-Surgery is built on a modular architecture with four decoupled registries---\texttt{MODEL}, \texttt{DATASET}, \texttt{TRAINER}, and \texttt{EVALUATOR}---orchestrated by Hydra configuration.
Figure~\ref{fig:architecture} illustrates the overall design.

\subsection{Unified Entry Point}

All three operations share a single training script (\texttt{train.py}) that dispatches to the appropriate pipeline based on the \texttt{mode} parameter:

\begin{lstlisting}
# Knowledge Editing
python src/train.py mode=edit \
  trainer=edit/ROME model=Llama-3.1-8B \
  experiment=edit/zsre/default

# Machine Unlearning
python src/train.py mode=unlearn \
  trainer=SimNPO model=Llama-3.1-8B \
  experiment=unlearn/tofu/default

# Knowledge Injection
python src/train.py mode=inject \
  trainer=inject/LoRA model=Qwen2.5-7B \
  experiment=inject/alpaca/default
\end{lstlisting}

\subsection{Registry-Based Architecture}

Each component type has a dedicated registry that maps string names to Python classes:

\begin{itemize}
\setlength\itemsep{0pt}
\item \texttt{TRAINER\_REGISTRY}: 48 method implementations (15 edit + 14 unlearn + 6 inject + 13 multimodal)
\item \texttt{DATASET\_REGISTRY}: 20+ dataset handlers covering QA, editing, injection, and multimodal formats
\item \texttt{MODEL\_REGISTRY}: supports any HuggingFace \texttt{AutoModelForCausalLM}-compatible model via YAML config
\item \texttt{EVALUATOR\_REGISTRY}: 18 evaluators spanning TOFU, MUSE, WMDP, and edit-specific metrics
\end{itemize}

Components are decoupled: methods do not depend on specific datasets or models.
A new component requires only (1) a Python class inheriting the base, (2) one registration line, and (3) a YAML config file.

\subsection{Three Knowledge Operation Pipelines}

\paragraph{Editing Pipeline.}
Edit requests pass through \texttt{build\_edit\_requests} $\to$ \texttt{execute\_edit\_requests} $\to$ \texttt{save\_edit\_artifacts}.
Each editor implements an \texttt{edit(requests)} method that modifies model weights and returns a backup for rollback.
Supported method families include locate-then-edit (ROME, MEMIT, AlphaEdit), meta-learning (MEND, MALMEN, InstructEdit), memory-based (GRACE, WISE, IKE, SERAC), and unstructured editing (UNKE, AnyEdit).

\paragraph{Unlearning Pipeline.}
Forget and retain datasets are combined into a \texttt{ForgetRetainDataset} that yields paired samples.
Each unlearning trainer overrides \texttt{compute\_loss} to implement its objective (gradient ascent, preference optimization, representation manipulation, etc.).

\paragraph{Injection Pipeline.}
PEFT adapters (LoRA, DoRA, AdaLoRA, LoReFT, BREP) are applied to the base model.
The injection trainer handles adapter initialization, training, and optional merging.

\begin{figure}[t]
\centering
\fbox{\parbox{0.95\columnwidth}{\centering\textit{[Architecture diagram placeholder]}\\
\small train.py $\to$ Registry dispatch $\to$ Edit/Unlearn/Inject pipelines\\
Hydra configs combine Model + Trainer + Data + Eval}}
\caption{Know-Surgery architecture. A unified entry point dispatches to three pipelines through registry-based component resolution.}
\label{fig:architecture}
\end{figure}

\section{Key Features}
\label{sec:features}

\subsection{Comprehensive Method Library}

Table~\ref{tab:methods} summarizes the methods integrated in Know-Surgery across three operation categories, totaling 35 text-based methods plus 13 multimodal methods.

\begin{table}[t]
\centering
\small
\setlength{\tabcolsep}{2pt}
\begin{tabular}{@{}llp{3.8cm}@{}}
\toprule
\textbf{Op.} & \textbf{Category} & \textbf{Methods} \\
\midrule
\multirow{5}{*}{\rotatebox{90}{Edit (15)}}
 & Locate-Edit & ROME, MEMIT, MEMIT\_M, AlphaEdit \\
 & Memory & GRACE, WISE, IKE, SERAC \\
 & Meta-learn & MEND, MALMEN, InstructEdit \\
 & Unified & UniKE, NMKE \\
 & Unstruct. & UNKE, AnyEdit \\
\midrule
\multirow{4}{*}{\rotatebox{90}{Unl. (14)}}
 & Gradient & GradAscent, GradDiff, SGA \\
 & Preference & NPO, DPO, SimNPO \\
 & Represent. & RMU, FLAT, SOUL \\
 & Others & UNDIAL, CEU, SatImp, WGA, PDU \\
\midrule
Inj. (6) & PEFT & LoRA, DoRA, AdaLoRA, LoReFT, BREP \\
\midrule
\multirow{2}{*}{MM (13)}
 & MM Unl. & MMGradAsc., MMNPO, MMVKD, ... (7) \\
 & MM Edit & MM-IKE, MM-UniKE, ... (6) \\
\bottomrule
\end{tabular}
\caption{Methods supported in Know-Surgery. MEMIT\_M = MEMIT\_Merge.}
\label{tab:methods}
\end{table}

\subsection{Multi-type Knowledge Coverage}

Unlike existing tools that focus primarily on factual triplets, Know-Surgery covers five knowledge types through 8+ editing datasets:

\begin{table}[t]
\centering
\small
\setlength{\tabcolsep}{2pt}
\begin{tabular}{@{}lp{2.6cm}p{2.8cm}@{}}
\toprule
\textbf{Type} & \textbf{Datasets} & \textbf{Methods} \\
\midrule
Factual & CF, ZsRE, LEME & ROME, MEMIT, AlphaEdit \\
Event & ELKEN, UnKE, AKEW & MEMIT, AnyEdit \\
Concept & ConceptEdit & ROME, AlphaEdit \\
Proced. & EditEvery & AnyEdit, UNKE \\
Rule & RuleTaker, FOLIO & BREP, LoReFT \\
\bottomrule
\end{tabular}
\caption{Knowledge type coverage. CF=CounterFact.}
\label{tab:knowledge}
\end{table}

\subsection{Multimodal Extension}

Know-Surgery extends to multimodal scenarios through two pathways:
\begin{itemize}
\setlength\itemsep{0pt}
\item \textbf{MM Unlearning} (7 methods): MMGradAscent, MMGradDiff, MMKLMin, MMNPO, MMRetainFT, MMUnlearner, MMVKD, evaluated on CLEAR and MLLMU-Bench.
\item \textbf{MM Editing} (6 methods): MM-IKE, MM-GRACE, MM-WISE, MM-MEND, MM-SERAC, MM-UniKE, extending text editing methods to vision-language models (InternVL, LLaVA, Qwen-VL) via a shared \texttt{MMEditMixin} base class.
\end{itemize}

\subsection{Three-Step Extension Protocol}

Adding a new component requires only three steps:
\begin{enumerate}
\setlength\itemsep{0pt}
\item \textbf{Implement}: Write a Python class inheriting the appropriate base class (\texttt{EditTrainer}, \texttt{FinetuneTrainer}, or \texttt{InjectTrainer}).
\item \textbf{Register}: Add one line to the registry (\texttt{\_register\_trainer(MyMethod)}).
\item \textbf{Configure}: Create a YAML file specifying the handler name and method-specific hyperparameters.
\end{enumerate}
No changes to the core framework, data loading, or evaluation code are needed.

\section{Demonstration}
\label{sec:demo}

Know-Surgery provides two interfaces for interactive exploration.

\paragraph{Command-Line Interface.}
Users run experiments with a single command, combining model, method, and dataset through Hydra overrides:

\begin{lstlisting}
# Edit: ROME on CounterFact with Qwen2.5-7B
python src/train.py mode=edit \
  trainer=edit/ROME model=Qwen2.5-7B-Instruct \
  experiment=edit/counterfact/default

# Unlearn: SimNPO on TOFU
python src/train.py mode=unlearn \
  trainer=SimNPO model=Llama-3.2-1B-Instruct \
  experiment=unlearn/tofu/default \
  task_name=my_simnpo_run

# Inject: LoRA on Alpaca
python src/train.py mode=inject \
  trainer=inject/LoRA model=Qwen2.5-7B-Instruct \
  experiment=inject/alpaca/default

# Evaluate a saved checkpoint
python src/eval.py \
  experiment=eval/tofu/default \
  model.model_args.pretrained_model_name_or_path=\
    saves/unlearn/my_simnpo_run
\end{lstlisting}

\paragraph{Web Interface.}
A Flask+Vue web UI (\texttt{new\_ui/}) allows users to select models, methods, and datasets from dropdown menus, configure hyperparameters, launch experiments, and compare results visually. The frontend is built with Vite + TypeScript and communicates with a Python Flask backend via REST APIs.

\begin{figure}[t]
\centering
\fbox{\parbox{0.95\columnwidth}{\centering\textit{[WebUI screenshot placeholder]}\\
\small Configuration panel + Results comparison view}}
\caption{Know-Surgery web interface for interactive experiment management.}
\label{fig:webui}
\end{figure}

\paragraph{Demo Video.}
A 2.5-minute screencast demonstrates the full workflow: selecting a method, configuring an experiment, running it, and analyzing results through both CLI and web interfaces.

\section{Evaluation}
\label{sec:eval}

We evaluate Know-Surgery on standard benchmarks to verify correctness and demonstrate cross-paradigm comparison capabilities.

\subsection{Knowledge Editing}

We compare representative editing methods on CounterFact and ZsRE. Metrics follow the standard four-dimensional evaluation: Reliability (edit success rate), Generalization (rephrase accuracy), Locality (unrelated knowledge preservation), and Portability (reasoning transfer).

\begin{table}[t]
\centering
\small
\begin{tabular}{lcccc}
\toprule
\textbf{Method} & \textbf{Rel.} & \textbf{Gen.} & \textbf{Loc.} & \textbf{Port.} \\
\midrule
\multicolumn{5}{c}{\textit{CounterFact}} \\
\midrule
ROME & -- & -- & -- & -- \\
MEMIT & -- & -- & -- & -- \\
MEND & -- & -- & -- & -- \\
AlphaEdit & -- & -- & -- & -- \\
GRACE & -- & -- & -- & -- \\
IKE & -- & -- & -- & -- \\
\midrule
\multicolumn{5}{c}{\textit{ZsRE}} \\
\midrule
ROME & -- & -- & -- & -- \\
MEMIT & -- & -- & -- & -- \\
MEND & -- & -- & -- & -- \\
AlphaEdit & -- & -- & -- & -- \\
GRACE & -- & -- & -- & -- \\
IKE & -- & -- & -- & -- \\
\bottomrule
\end{tabular}
\caption{Knowledge editing results. Rel.=Reliability, Gen.=Generalization, Loc.=Locality, Port.=Portability. Values are percentages. \textit{[TO BE FILLED after experiments]}}
\label{tab:edit}
\end{table}

\subsection{Machine Unlearning}

We evaluate unlearning methods on the TOFU benchmark (forget10 split). Forget Quality measures how well the model forgets target data; Model Utility measures retained performance on non-target data.

\begin{table}[t]
\centering
\small
\begin{tabular}{lcc}
\toprule
\textbf{Method} & \textbf{Forget Quality} & \textbf{Model Utility} \\
\midrule
GradAscent & -- & -- \\
GradDiff & -- & -- \\
NPO & -- & -- \\
DPO & -- & -- \\
SimNPO & -- & -- \\
RMU & -- & -- \\
CEU & -- & -- \\
PDU & -- & -- \\
WGA & -- & -- \\
\bottomrule
\end{tabular}
\caption{Machine unlearning results on TOFU (forget10). \textit{[TO BE FILLED]}}
\label{tab:unlearn}
\end{table}

\subsection{Comparison with Existing Toolkits}

\begin{table}[t]
\centering
\small
\setlength{\tabcolsep}{3pt}
\begin{tabular}{@{}lcccc@{}}
\toprule
\textbf{Feature} & \textbf{Ours} & \textbf{EasyEdit} & \textbf{OpenUnl.} \\
\midrule
Edit methods & 15 & 15+ & 0 \\
Unlearn methods & 14 & 0 & 11 \\
Inject methods & 6 & 1 & 5 \\
Knowledge types & 5 & 1--2 & N/A \\
Multimodal & \checkmark & \checkmark & \checkmark \\
Unified framework & \checkmark & $\times$ & $\times$ \\
Config system & Hydra & Py API & Hydra \\
Web UI & \checkmark & \checkmark & $\times$ \\
Model configs & 42 & 10+ & 10+ \\
\bottomrule
\end{tabular}
\caption{Comparison with existing toolkits.}
\label{tab:compare}
\end{table}

\section{Conclusion}
\label{sec:conclusion}

We presented Know-Surgery, a unified toolkit for controllable knowledge update in large language models. By integrating knowledge editing, machine unlearning, and knowledge injection within a single registry-based framework, Know-Surgery enables researchers to directly compare and combine these complementary paradigms. The toolkit supports 35 text-based methods and 13 multimodal methods across 42 model configurations, covers 5 knowledge types, and is easily extensible through a three-step protocol. Know-Surgery is open-source and includes both CLI and web interfaces for accessibility.

\section*{Limitations}

Know-Surgery currently focuses on decoder-only language models. Support for encoder-decoder architectures (e.g., T5) is limited to methods that do not require specific layer path configurations. Evaluation of procedural and rule-based knowledge editing lacks established benchmarks, and our coverage in these areas relies on injection-based methods rather than direct editing.

\section*{Ethics Statement}

Knowledge editing and unlearning tools carry inherent dual-use risks. While designed to improve model safety and factual accuracy, they could potentially be misused to inject misinformation or selectively remove safety-relevant knowledge. We encourage responsible use and recommend that practitioners document all knowledge modifications applied to deployed models.


\bibliography{references}

\appendix
\section{Complete Method Reference}
\label{sec:appendix-methods}

\begin{table*}[t]
\centering
\footnotesize
\setlength{\tabcolsep}{4pt}
\renewcommand{\arraystretch}{0.9}
\begin{minipage}[t]{0.49\textwidth}\centering
\begin{tabular}{@{}lll@{}}
\toprule
\textbf{Method} & \textbf{Category} & \textbf{Reference} \\
\midrule
\multicolumn{3}{@{}l}{\textit{Editing (15)}} \\
ROME & Locate-Edit & \citet{meng2022rome} \\
MEMIT & Locate-Edit & \citet{meng2023memit} \\
MEMIT\_Merge & Locate-Edit & \citet{meng2023memit} \\
AlphaEdit & Locate-Edit & \citet{fang2024alphaedit} \\
MEND & Meta-learn & \citet{mitchell2022mend} \\
MALMEN & Meta-learn & \citet{tan2024malmen} \\
InstructEdit & Meta-learn & \citet{tian2024instructedit} \\
GRACE & Memory & \citet{hartvigsen2023grace} \\
WISE & Memory & \citet{wang2024wise} \\
IKE & Memory & \citet{zheng2023ike} \\
SERAC & Memory & \citet{mitchell2022serac} \\
UniKE & Unified & \citet{pan2024unike} \\
NMKE & Unified & \citet{liu2025nmke} \\
UNKE & Unstruct. & \citet{wang2024unke} \\
AnyEdit & Multi-strat. & \citet{wang2024unke}$^\dagger$ \\
\midrule
\multicolumn{3}{@{}l}{\textit{Injection (6)}} \\
LoRA & PEFT & \citet{hu2022lora} \\
DoRA & PEFT & \citet{liu2024dora} \\
AdaLoRA & PEFT & \citet{zhang2023adalora} \\
LoReFT & PEFT & \citet{wu2024loreft} \\
BREP & PEFT & \citet{liang2025brep} \\
\bottomrule
\end{tabular}
\end{minipage}\hfill
\begin{minipage}[t]{0.49\textwidth}\centering
\begin{tabular}{@{}lll@{}}
\toprule
\textbf{Method} & \textbf{Category} & \textbf{Reference} \\
\midrule
\multicolumn{3}{@{}l}{\textit{Unlearning (14)}} \\
GradAscent & Gradient & \citet{yao2023unlearning} \\
GradDiff & Gradient & \citet{maini2024tofu} \\
SGA & Gradient & \citet{pang2026sga} \\
NPO & Preference & \citet{zhang2024npo} \\
SimNPO & Preference & \citet{fan2024simnpo} \\
DPO & Preference & \citet{rafailov2023dpo} \\
RMU & Represent. & \citet{li2024wmdp} \\
FLAT & Represent. & \citet{zhang2025flat} \\
SOUL & Represent. & \citet{jia2024soul} \\
UNDIAL & Other & \citet{dong2025undial} \\
CEU & Other & \citet{yang2025ceu} \\
SatImp & Other & \citet{yang2025satimp} \\
WGA & Other & \citet{geffect2025wga} \\
PDU & Other & \citet{entesari2025pdu} \\
\midrule
\multicolumn{3}{@{}l}{\textit{Multimodal (13)}} \\
MMGradAscent & MM-Unl. & \citet{yao2023unlearning} \\
MMGradDiff & MM-Unl. & \citet{maini2024tofu} \\
MMKLMin & MM-Unl. & -- \\
MMNPO & MM-Unl. & \citet{zhang2024npo} \\
MMRetainFT & MM-Unl. & -- \\
MMUnlearner & MM-Unl. & \citet{huo2025mmunlearner} \\
MMVKD & MM-Unl. & \citet{wang2024mmvkd} \\
MM-IKE & MM-Edit & \citet{zheng2023ike} \\
MM-GRACE & MM-Edit & \citet{hartvigsen2023grace} \\
MM-WISE & MM-Edit & \citet{wang2024wise} \\
MM-MEND & MM-Edit & \citet{mitchell2022mend} \\
MM-SERAC & MM-Edit & \citet{mitchell2022serac} \\
MM-UniKE & MM-Edit & \citet{pan2024unike} \\
\bottomrule
\end{tabular}
\end{minipage}
\caption{Complete method reference. 48 methods total: 15 text edit + 14 text unlearn + 6 inject (5 PEFT + full fine-tuning) + 7 MM unlearn + 6 MM edit. $^\dagger$This AnyEdit aggregates MEMIT, AlphaEdit, and UNKE strategies in a unified interface (distinct from the same-named ICML'25 method).}
\label{tab:all-methods}
\end{table*}

\section{Dataset Summary}
\label{sec:appendix-data}

\begin{table}[t]
\centering
\footnotesize
\renewcommand{\arraystretch}{0.9}
\begin{tabular}{llr}
\toprule
\textbf{Dataset} & \textbf{Knowledge Type} & \textbf{Size} \\
\midrule
CounterFact & Factual & 21,919 \\
ZsRE & Factual & 10,000 \\
LEME & Factual (long-form) & 7,390 \\
ELKEN & Event & 1,847 \\
UnKE & Event/Temporal & 1,000 \\
AKEW & Event/Multi-hop & 975+ \\
ConceptEdit & Conceptual & 452 \\
EditEvery & Procedural & 552 \\
MQuAKE & Multi-hop & 3,000+ \\
\midrule
TOFU & Unlearn & 4,000 \\
MUSE-News & Unlearn & 6.5M tokens \\
MUSE-Books & Unlearn & 6.5M tokens \\
WMDP & Unlearn (safety) & 3,668 \\
RWKU & Unlearn & 200 entities \\
\midrule
CLEAR & MM Unlearn & 3,700 \\
MLLMU-Bench & MM Unlearn & 500+ \\
MMEdit & MM Edit & 1,000+ \\
MMKE-Bench & MM Edit & 2,000+ \\
VLKEB & MM Edit & 1,500+ \\
\midrule
Alpaca & Inject & 52,000 \\
ShareGPT & Inject & 90,000+ \\
\bottomrule
\end{tabular}
\caption{Datasets supported in Know-Surgery.}
\label{tab:all-data}
\end{table}

\section{Supported Models}
\label{sec:appendix-models}

Know-Surgery provides 42 pre-configured model YAML files covering:
\begin{itemize}
\setlength\itemsep{0pt}
\item \textbf{Llama family}: Llama-2-7B/13B, Llama-3.1-8B, Llama-3.2-1B/3B
\item \textbf{Qwen family}: Qwen2.5-0.5B/1.5B/3B/7B/14B-Instruct, Qwen2.5VL-7B, Qwen2VL-2B, Qwen3-VL-4B
\item \textbf{Others}: Mistral-7B, Phi-3/3.5, Gemma-2B/7B/9B, Yi-1.5-6B/9B, GLM-4-9B, InternLM2.5-7B, DeepSeek-R1-Distill, GPT-2/J, Baichuan2-7B/13B
\item \textbf{Vision-Language}: InternVL2.5-2B/8B, LLaVA-1.5-7B/13B, LLaVA-OneVision-7B, InstructBLIP-Vicuna-7B, BLIP2-OPT-2.7B
\end{itemize}
Any HuggingFace-compatible model can be added by creating a single YAML configuration file.


\end{document}
